工业过程通常具有长时滞、变量耦合以及直接试错成本高等特点。仅依靠历史轨迹训练策略，难以在执行前回答“某个控制动作将使设备状态如何演化”；直接在真实设备上反复搜索动作则可能带来能耗、磨损和安全风险。本项目从已有工业轨迹学习可递归预测未来状态的世界模型，再利用冻结模型进行动作序列规划或策略学习，最后在 NeoRL 仿真环境中验证长期控制效果。

整体技术路线见\cref{fig:framework}。历史状态序列和动作首先用于训练世界模型；模型通过单步预测和多步递归预测评价后被冻结，分别支持基于模型的规划方法 CEM、iCEM、MPPI，以及基于模型的强化学习方法 MB-PPO。经过优化的策略最终在 NeoRL 仿真环境中执行，以1000步累计奖励衡量长期控制效果。

\begin{figure}[H]
  \centering
  \includegraphics[width=\textwidth]{fig1_framework.pdf}
  \caption{\textbf{项目整体技术路线。} NeoRL 历史轨迹被整理为历史状态序列与连续动作，用于训练下一状态预测模型。世界模型经过单步与多步预测评价后被冻结，并分别服务于基于模型的规划和基于模型的强化学习；优化后的控制策略在 NeoRL 仿真环境中验证长期累计奖励。}
  \label{fig:framework}
\end{figure}

本项目比较五类世界模型、三种规划方法和一种基于模型的强化学习方法。最终在 M100 上选择 GRU+CEM，在 M1000 与 M10000 上选择 Transformer-2L+MPPI。正文依次介绍状态与奖励定义、世界模型构建与评价、策略优化方法、仿真实验结果，以及这些结果对工业部署的启示。
